\chapter{Introduzione}
\label{cap:01-introduzione}

\section{Il contesto: che cosa promette il modello serverless}

Il modello \emph{Function-as-a-Service} (FaaS) propone al programmatore un
contratto tanto semplice quanto ambizioso: si consegna alla piattaforma una
funzione, ossia un frammento di codice senza stato con una firma nota, e la
piattaforma si assume ogni responsabilit\`a residua --- collocarla su una
macchina, replicarla quando gli arrivi crescono, ritirarla quando cessano,
instradare le richieste, riprovare i fallimenti, contabilizzare l'uso. Il
programmatore non dimensiona nulla e non paga nulla quando nessuno invoca.

\`E un contratto attraente perch\'e sposta un intero registro di decisioni
--- capacit\`a, elasticit\`a, collocazione --- dall'utente al fornitore. Ed \`e
anche un contratto costoso, perch\'e quelle decisioni non spariscono: si
concentrano dentro la piattaforma, dove diventano politiche automatiche che
l'utente non vede e, quando il comportamento osservato non lo soddisfa, non sa
come correggere. Il \emph{cold start} \`e l'esempio pi\`u noto di questa
opacit\`a: un ritardo che il chiamante misura ma che dipende da una decisione
di collocazione presa altrove, in un istante che il chiamante non conosce, sulla
base di uno stato che non gli viene esposto.

La conseguenza per chi fa ricerca su queste piattaforme \`e concreta. Studiare
una politica di scaling, di ammissione o di controllo della concorrenza richiede
di poterla sostituire e di poter attribuire una differenza misurata a quella
sostituzione e non ad altro. Sulle piattaforme FaaS mature questo \`e difficile
per una ragione strutturale, non accidentale: sono sistemi distribuiti composti
da molti processi cooperanti --- un ingress, un autoscaler, un attivatore, un
controller, una coda, un data plane --- ciascuno con la propria dinamica, il
proprio periodo di campionamento e le proprie isteresi. Il percorso di una
richiesta attraversa diversi confini di processo prima di raggiungere il codice
dell'utente, e ogni confine contribuisce alla latenza in un modo che
l'osservatore non pu\`o disaggregare senza strumentare l'intero sistema.

\section{Il problema}

Da qui la domanda che motiva questo lavoro: \emph{qual \`e la piattaforma pi\`u
piccola che \`e ancora una piattaforma FaaS, e che cosa diventa misurabile una
volta che la si \`e ridotta a quella dimensione?}

La domanda non \`e retorica. Una piattaforma pu\`o essere ridotta lungo molte
dimensioni, e non tutte le riduzioni sono equivalenti. Togliere
l'autenticazione riduce il codice ma non accorcia il percorso critico di
un'invocazione; togliere la persistenza dello stato lo accorcia ma rinuncia alla
sopravvivenza a un riavvio; togliere l'alta disponibilit\`a elimina un intero
registro di problemi di consenso ma vincola la piattaforma a un singolo punto di
fallimento. La tesi di questo lavoro \`e che esiste un insieme di rinunce
--- pod singolo, stato in memoria, nessuna autenticazione --- che, prese
insieme, producono un sistema il cui percorso critico \`e leggibile per intero
da una sola persona e in cui una modifica di politica \`e osservabile senza
rumore di fondo, e che tale sistema resta abbastanza completo da eseguire
carichi reali su Kubernetes.

\nf{} \`e la realizzazione di questa tesi. \`E un control plane scritto in Java
su Spring WebFlux, il cui nucleo conta circa 5\,000 righe di codice di
produzione, esteso da nove moduli opzionali che ne aggiungono complessivamente
circa 6\,800. Il rapporto fra le due cifre \`e esso stesso una scelta di
progetto: i moduli, presi insieme, sono pi\`u grandi del nucleo che estendono,
il che significa che la maggior parte di ci\`o che rende \nf{} interessante \`e
anche ci\`o che si pu\`o togliere senza che smetta di funzionare.

\section{Domande di ricerca}

Il lavoro \`e organizzato attorno a quattro domande.

\begin{description}[style=nextline, leftmargin=0pt]
  \item[\textbf{RQ1 --- Modularit\`a a costo zero sul percorso critico.}]
  \`E possibile costruire un control plane in cui funzionalit\`a come le code,
  l'autoscaling e il controllo della concorrenza siano \emph{assenti} anzich\'e
  disattivate, senza che la loro assenza richieda rami condizionali sul percorso
  di invocazione? Il \cref{cap:08-sistema-moduli} descrive la soluzione adottata
  --- selezione a tempo di build attraverso un SPI, con implementazioni no-op di
  default per ogni punto di estensione --- e ne discute i limiti.

  \item[\textbf{RQ2 --- Da quale segnale deve decidere un controllore di
  concorrenza.}] Il limite di concorrenza per funzione \`e la grandezza che
  determina simultaneamente l'utilizzazione delle repliche e il tempo che un
  chiamante passa in attesa. La domanda \`e da quale segnale misurato tale limite
  debba essere derivato. Il \cref{cap:12-concurrency-control} mostra
  analiticamente e sperimentalmente che il tempo di servizio --- la grandezza
  che il control plane misura pi\`u naturalmente, fra dispatch e completamento
  --- \`e la scelta sbagliata, e che il tempo di sojourn \`e quella corretta.

  \item[\textbf{RQ3 --- Concorrenza come risorsa condivisa.}] Quando pi\`u
  funzioni insistono sulle stesse repliche o sullo stesso nodo, un controllore
  che chiede a ciascuna funzione di trovare da sola il proprio limite le sta
  chiedendo di decidere su una risorsa condivisa senza informazioni sui vicini.
  Il \cref{cap:12-concurrency-control} formula un'alternativa (modo
  \code{BUDGETED}) che separa la stima del fabbisogno dall'allocazione, e il
  \cref{cap:25-valutazione-concorrenza} riporta la misura del cross-talk fra
  funzioni co-residenti.

  \item[\textbf{RQ4 --- Qual \`e il costo reale di un control plane.}] Un
  processo che non esegue codice utente ma lo orchestra ha un'impronta di
  memoria e un profilo di avvio propri. Il \cref{cap:24-footprint} misura dove
  va la memoria di una JVM Spring, che cosa cambia compilando lo stesso codice
  in immagine nativa GraalVM, e come la scelta del garbage collector interagisce
  con i limiti di memoria imposti dal container.
\end{description}

\section{Contributi}

Il documento presenta i seguenti contributi.

\begin{enumerate}[leftmargin=*]
  \item \textbf{Un'architettura di control plane FaaS a nucleo minimale e moduli
  opzionali}, con un meccanismo di composizione a tempo di build che consente di
  produrre binari privi delle funzionalit\`a non selezionate, e un insieme di
  punti di estensione i cui default no-op degradano il sistema in modo esplicito
  anzich\'e farlo fallire (\cref{cap:04-architettura-control-plane,cap:06-nucleo,cap:08-sistema-moduli}).

  \item \textbf{Una tassonomia operativa di cinque politiche di controllo della
  concorrenza per funzione} --- fissa, statica per replica, adattiva per replica,
  a budget di piattaforma, e a minimizzazione del tempo di sojourn --- con la
  discussione analitica del perch\'e le prime quattro non possono essere valutate
  con lo stesso esperimento (\cref{cap:12-concurrency-control}).

  \item \textbf{Due risultati sperimentali negativi}, riportati come tali: un
  controllore che decide dal tempo di servizio non ammette ottimo interno e
  converge al proprio limite inferiore; e un generatore di carico a ciclo chiuso
  non pu\`o distinguere fra le politiche, perch\'e retroaziona sulla stessa
  grandezza che il controllore sta manipolando
  (\cref{cap:25-valutazione-concorrenza}).

  \item \textbf{Un modulo di offload trasparente edge$\to$cloud} attivabile in
  modo eager per funzione o reattivo sulla pressione della coda sincrona, con una
  semantica a singolo salto e senza fallback locale, e la giustificazione di
  entrambe le rinunce (\cref{cap:13-offload}).

  \item \textbf{Una caratterizzazione dell'impronta di runtime} del control plane
  su JVM e in immagine nativa, inclusa l'interazione fra scelta del collector e
  limite di memoria del container (\cref{cap:24-footprint}).

  \item \textbf{Un impianto sperimentale riproducibile}, esterno al codice della
  piattaforma, che provvisiona l'ambiente, esegue i carichi e confronta i
  risultati contro record versionati su profilo hardware fissato
  (\cref{cap:22-metodo-sperimentale}).
\end{enumerate}

\section{Che cosa questo lavoro non \`e}

\`E utile dichiarare subito i confini, per non farli scoprire al lettore a
met\`a documento.

\nf{} non compete con Knative, OpenWhisk o OpenFaaS su completezza funzionale, e
non \`e proposto come loro sostituto in produzione. Non implementa
autenticazione n\'e autorizzazione, quindi non \`e esponibile su una rete non
fidata. Non replica il proprio stato: un riavvio del pod perde il registro delle
funzioni e lo storico delle esecuzioni. Non implementa alta disponibilit\`a, e
l'esistenza di un singolo processo di controllo \`e assunta ovunque nel codice,
non incapsulata dietro un'astrazione che un giorno potrebbe diventare
distribuita.

Sul piano sperimentale, i numeri riportati nel \cref{cap:23-prestazioni-base} e
seguenti sono ottenuti su un profilo hardware specifico e dichiarato. Non sono
proposti come caratterizzazione assoluta della piattaforma n\'e come base per un
confronto con altre piattaforme misurate altrove. Il
\cref{cap:27-discussione} enumera in modo esplicito le affermazioni che i dati
raccolti non sostengono.

\section{Struttura del documento}

Il documento \`e diviso in sei parti.

La \textbf{Parte I} (\crefrange{cap:01-introduzione}{cap:03-requisiti-principi})
inquadra il problema, rivede lo stato dell'arte e formalizza i vincoli di
progetto, mostrando per ciascuno quale esperimento abilita.

La \textbf{Parte II} (\crefrange{cap:04-architettura-control-plane}{cap:08-sistema-moduli})
descrive l'architettura: la vista d'insieme, i contratti di dati e l'API REST,
il nucleo e i suoi punti di estensione, le modalit\`a di esecuzione e il
meccanismo di composizione modulare.

La \textbf{Parte III} (\crefrange{cap:09-async-queue}{cap:16-moduli-diagnostici})
dedica un capitolo a ciascun modulo opzionale. Il
\cref{cap:12-concurrency-control} \`e il pi\`u lungo della parte ed \`e quello a
contenuto scientifico maggiore.

La \textbf{Parte IV} (\crefrange{cap:17-runtime-funzione}{cap:20-build-packaging})
attraversa il lato funzione: il contratto di runtime, gli SDK nei vari
linguaggi, gli strumenti a riga di comando e la catena di build.

La \textbf{Parte V} (\crefrange{cap:21-osservabilita}{cap:25-valutazione-concorrenza})
presenta il modello di osservabilit\`a, il metodo sperimentale e le tre
campagne di misura.

La \textbf{Parte VI} (\crefrange{cap:26-qualita-verifica}{cap:29-conclusioni})
discute la verifica del software, i limiti dei risultati e le direzioni
successive.

\section{Convenzioni}

Il codice sorgente \`e citato con il percorso relativo alla radice del
repository, per esempio \code{platform/control-plane/src/main/java/\ldots}. Il
package radice del codice Java \`e \code{it.unimib.datai.nanofaas}; nel testo i
nomi di classe sono riportati senza package quando non ambigui. I nomi dei
moduli opzionali compaiono nella forma con cui sono selezionabili dalla riga di
comando (\code{async-queue}, \code{concurrency-control}, \ldots), che coincide
con il nome della directory sotto \code{platform/modules/}.

Le grandezze temporali sono espresse in millisecondi salvo indicazione contraria.
Con \emph{tempo di servizio} si intende l'intervallo fra l'istante in cui il
control plane consegna l'invocazione al runtime e l'istante in cui ne riceve la
risposta; con \emph{tempo di attesa} l'intervallo fra l'ammissione della
richiesta e la consegna; con \emph{tempo di sojourn} la loro somma, che \`e ci\`o
che il chiamante osserva al netto della rete. La distinzione \`e centrale nel
\cref{cap:12-concurrency-control} e viene mantenuta rigorosamente in tutto il
documento.

Infine, i riquadri marcati \textbf{Da misurare} segnalano affermazioni per cui
il progetto dispone di uno strumento di misura ma non ancora di un dato
raccolto. Sono lasciati visibili di proposito: un documento che presenta una
piattaforma come strumento sperimentale deve distinguere ci\`o che ha misurato
da ci\`o che potrebbe misurare.
